\documentclass[../../main.tex]{subfiles}% 注意这里的文件路径不能用 ./main.tex ,否则用latexmk编译子文件会报错
\graphicspath{{\subfix{./image/}}{\subfix{../../image/}}} % 指定图片引用路径,后续可以直接使用图片文件名
% 如果图片存放在上级目录,则使用 ../../image/ 作为备用路径

% 例如:
% \begin{figure}[H]
% \centering
% \includegraphics[width=0.3\textwidth]{图.png}
% \caption{}
% \label{figure:图}
% \end{figure}
% 注意:上述\label{}一定要放在\caption{}之后,否则引用图片序号会只会显示??.

\begin{document}

\section{集族}\label{section:ch01-集族}

\subsection{可测空间}\label{subsection:ch01-可测空间}

\begin{definition}\label{definition:ch01-半代数代数}
设 $X$ 为一集合，非空集族 $\mathcal{S}\subseteq\mathcal{P}(X)$ 称为一个\textbf{半代数}，若：
\begin{enumerate}[(i)]
\item $E,F\in\mathcal{S}\Rightarrow E\cap F\in\mathcal{S}$；
\item 对任意 $E\in\mathcal{S}$，存在 $E^c$ 的一个在$\mathcal{S}$中的划分。
\end{enumerate}
集族 $\mathcal{A}\subseteq\mathcal{P}(X)$ 称为一个\textbf{代数}，若：
\begin{enumerate}[(i)]
\item $\varnothing \in\mathcal{A}$；
\item $E\in\mathcal{A}\Rightarrow E^c\in\mathcal{A}$；
\item $E,F\in\mathcal{A}\Rightarrow E\cup F\in\mathcal{A}$。
\end{enumerate}
\end{definition}
\begin{remark}
显然$\mathcal{P}(X)$既是半代数，也是代数.
\end{remark}

\begin{proposition}
对任意集族 $\mathcal{F}\subseteq\mathcal{P}(X)$，交集
\[
\mathcal{A}_0(\mathcal{F})\triangleq \bigcap\{\mathcal{A}\mid\mathcal{A}\text{ 为代数，}\mathcal{A}\supseteq\mathcal{F}\}
\]
是一个代数，它包含 $\mathcal{F}$ 且被包含在所有包含 $\mathcal{F}$ 的代数 $\mathcal{A}$ 中。它称为\textbf{由 $\mathcal{F}$ 生成的代数}，或\textbf{包含 $\mathcal{F}$ 的最小代数}。
\end{proposition}
\begin{proof}
\noindent$\mathbf{Step}\mathbf{1}.${\heiti 说明所有包含 $\mathcal{F}$ 的代数构成的集族非空.}
记 $\mathfrak{A}$ 为所有包含 $\mathcal{F}$ 的代数组成的集族, 则 $\mathcal{A}_0(\mathcal{F})=\bigcap\limits_{\mathcal{A}\in\mathfrak{A}}\mathcal{A}$。由于 $\mathcal{P}(X)$ 是代数且 $\mathcal{P}(X)\supseteq\mathcal{F}$, 故 $\mathcal{P}(X)\in\mathfrak{A}$, 从而 $\mathfrak{A}\neq\varnothing$。因此 $\mathcal{A}_0(\mathcal{F})$ 的定义合理。

\noindent$\mathbf{Step}\mathbf{2}.${\heiti 说明 $\mathcal{F}\subseteq\mathcal{A}_0(\mathcal{F})$.}
任取 $E\in\mathcal{F}$。对任意 $\mathcal{A}\in\mathfrak{A}$, 由 $\mathcal{A}\supseteq\mathcal{F}$ 得 $E\in\mathcal{A}$。于是 $E\in\bigcap\limits_{\mathcal{A}\in\mathfrak{A}}\mathcal{A}=\mathcal{A}_0(\mathcal{F})$。故 $\mathcal{F}\subseteq\mathcal{A}_0(\mathcal{F})$。

\noindent$\mathbf{Step}\mathbf{3}.${\heiti 验证 $\varnothing\in\mathcal{A}_0(\mathcal{F})$.}
对任意 $\mathcal{A}\in\mathfrak{A}$, 按\cref{definition:ch01-半代数代数}中代数的定义有 $\varnothing\in\mathcal{A}$。故 $\varnothing\in\bigcap\limits_{\mathcal{A}\in\mathfrak{A}}\mathcal{A}=\mathcal{A}_0(\mathcal{F})$。

\noindent$\mathbf{Step}\mathbf{4}.${\heiti 验证 $\mathcal{A}_0(\mathcal{F})$ 对补运算封闭.}
任取 $E\in\mathcal{A}_0(\mathcal{F})$。对任意 $\mathcal{A}\in\mathfrak{A}$, 有 $E\in\mathcal{A}$。因 $\mathcal{A}$ 是代数, 故 $E^c\in\mathcal{A}$。由 $\mathcal{A}$ 的任意性得 $E^c\in\mathcal{A}_0(\mathcal{F})$。因此 $\mathcal{A}_0(\mathcal{F})$ 对补运算封闭。

\noindent$\mathbf{Step}\mathbf{5}.${\heiti 验证 $\mathcal{A}_0(\mathcal{F})$ 对有限并封闭并确认其为代数.}
任取 $E,F\in\mathcal{A}_0(\mathcal{F})$。对任意 $\mathcal{A}\in\mathfrak{A}$, 有 $E,F\in\mathcal{A}$。因 $\mathcal{A}$ 是代数, 故 $E\cup F\in\mathcal{A}$。由 $\mathcal{A}$ 的任意性得 $E\cup F\in\mathcal{A}_0(\mathcal{F})$。因此 $\mathcal{A}_0(\mathcal{F})$ 对有限并封闭。结合 $\mathbf{Step3}$ 至 $\mathbf{Step5}$ 及\cref{definition:ch01-半代数代数}中代数定义可知 $\mathcal{A}_0(\mathcal{F})$ 是代数。

\noindent$\mathbf{Step}\mathbf{6}.${\heiti 说明最小性.}
设 $\mathcal{B}$ 是代数且 $\mathcal{B}\supseteq\mathcal{F}$。则 $\mathcal{B}\in\mathfrak{A}$, 于是 $\mathcal{A}_0(\mathcal{F})=\bigcap\limits_{\mathcal{A}\in\mathfrak{A}}\mathcal{A}\subseteq\mathcal{B}$。故 $\mathcal{A}_0(\mathcal{F})$ 被包含在所有包含 $\mathcal{F}$ 的代数中。
\end{proof}

\begin{example}\label{example:ch01-N胞腔}
设 $a\equiv(a_1,\cdots ,a_n)$、$b\equiv(b_1,\cdots ,b_n)$，其中 $a_i,b_i\in[-\infty,\infty]$ 且对每个 $i=1,\cdots ,n$（$n\in\mathbb{N}$）有 $a_i<b_i$。考虑 $N$--胞腔
\[
(a,b]\triangleq \{x\equiv(x_1,\cdots ,x_n)\mid x_i\in[-\infty,\infty],\ a_i<x_i\leqslant b_i\ \forall i=1,\cdots ,n\}
\]
其体积为
\[
\operatorname{vol}((a,b])\triangleq \prod_{i=1}^n(b_i-a_i)\leqslant\infty.
\]
容易看出，$N$--胞腔族 $\mathcal{I}_n$ 是一个半代数，且由 $\mathcal{I}_n$ 生成的代数 $\mathcal{A}_0(\mathcal{I}_n)$ 由 $N$--胞腔的有限不交并组成，即对每个 $E\in\mathcal{A}_0(\mathcal{I}_n)$，存在不交的 $(a_1,b_1],\cdots ,(a_p,b_p]$ 使得
\[
E=\bigcup_{k=1}^p(a_k,b_k].
\]
\end{example}

\begin{definition}\label{definition:ch01-sigma代数}
族 $\mathcal{A}\subseteq\mathcal{P}(X)$ 称为一个 $\sigma$--代数，若：
\begin{enumerate}[(i)]
\item $\varnothing \in\mathcal{A}$；
\item $E\in\mathcal{A}\Rightarrow E^c\in\mathcal{A}$；
\item 对任意序列 $\{E_k\}\subseteq\mathcal{A}$，有 $\bigcup_{k=1}^\infty E_k\in\mathcal{A}$。
\end{enumerate}
\end{definition}

\begin{definition}\label{definition:ch01-可测空间}
设 $\mathcal{A}\subseteq\mathcal{P}(X)$ 为一 $\sigma$--代数。偶 $(X,\mathcal{A})$ 称为一个可测空间，$\mathcal{A}$ 的元素称为可测集。
\end{definition}

\begin{remark}\label{remark:ch01-迹}
设 $F\subseteq X$。容易看出，若 $\mathcal{A}$ 为一 $\sigma$--代数，则族
\[
\mathcal{A}\cap F\triangleq \{E\cap F\mid E\in\mathcal{A}\}\subseteq\mathcal{P}(F)
\]
也是一个 $\sigma$--代数，称为 $\mathcal{A}$ 在 $F$ 上的迹。可测空间 $(F,\mathcal{A}\cap F)$ 称为 $(X,\mathcal{A})$ 的一个可测子空间。
\end{remark}

\begin{definition}\label{definition:ch01-单调类}
非空族 $\mathcal{M}\subseteq\mathcal{P}(X)$ 称为一个单调类，若对任意非减（相应地，非增）序列 $\{E_k\}\subseteq\mathcal{M}$，有 $\bigcup_{k=1}^\infty E_k\in\mathcal{M}$（相应地，$\bigcap_{k=1}^\infty E_k\in\mathcal{M}$）。
\end{definition}

显然，每个 $\sigma$--代数都是一个单调类。与代数的情形一样，对任意族 $\mathcal{F}\subseteq\mathcal{P}(X)$，我们可以考虑由 $\mathcal{F}$ 生成的最小 $\sigma$--代数
\[
\sigma_0(\mathcal{F})\triangleq \bigcap\{\mathcal{A}\mid\mathcal{A}\text{ 为 }\sigma\text{--代数，}\mathcal{A}\supseteq\mathcal{F}\}
\]
和最小单调类
\[
\mathcal{M}_0(\mathcal{F})\triangleq \bigcap\{\mathcal{M}\mid\mathcal{M}\text{ 为单调类，}\mathcal{M}\supseteq\mathcal{F}\}.
\]

\begin{definition}\label{definition:ch01-可分可测空间}\label{definition:ch01-可分空间}
可测空间 $(X,\mathcal{A})$ 称为可分的，若存在可数族 $\mathcal{S}\subseteq\mathcal{P}(X)$ 使得 $\mathcal{A}=\sigma_0(\mathcal{S})$。
\end{definition}

\begin{theorem}\label{theorem:ch01-单调类定理}
设 $\mathcal{A}\subseteq\mathcal{P}(X)$ 为一代数。则最小单调类 $\mathcal{M}_0(\mathcal{A})$ 与最小 $\sigma$--代数 $\sigma_0(\mathcal{A})$ 重合。
\end{theorem}

\begin{proof}
显然 $\varnothing \in\mathcal{M}_0(\mathcal{A})$。我们首先证明单调类 $\mathcal{M}_0(\mathcal{A})$ 是一个代数。令 $\widetilde{\mathcal{A}}\triangleq \{E\in\mathcal{M}_0(\mathcal{A})\mid E^c\in\mathcal{M}_0(\mathcal{A})\}$。由于 $\mathcal{A}$ 是代数且 $\mathcal{A}\subseteq\mathcal{M}_0(\mathcal{A})$，故 $\mathcal{A}\subseteq\widetilde{\mathcal{A}}$。另一方面，容易看出 $\widetilde{\mathcal{A}}$ 是一个单调类，因而 $\mathcal{M}_0(\mathcal{A})\subseteq\mathcal{M}_0(\widetilde{\mathcal{A}})\subseteq\widetilde{\mathcal{A}}$。由定义 $\widetilde{\mathcal{A}}\subseteq\mathcal{M}_0(\mathcal{A})$，故 $\widetilde{\mathcal{A}}=\mathcal{M}_0(\mathcal{A})$。因此对任意 $E\in\mathcal{M}_0(\mathcal{A})$，有 $E^c\in\mathcal{M}_0(\mathcal{A})$。类似地可证，对任意 $E,F\in\mathcal{M}_0(\mathcal{A})$，有 $E\cup F\in\mathcal{M}_0(\mathcal{A})$。故结论成立。

设 $\{E_k\}\subseteq\mathcal{M}_0(\mathcal{A})$，令 $F_j\triangleq \bigcup_{k=1}^j E_k$（$j,k\in\mathbb{N}$）。则序列 $\{F_j\}$ 非减，且 $\{F_j\}\subseteq\mathcal{M}_0(\mathcal{A})$，因为 $\mathcal{M}_0(\mathcal{A})$ 是一个代数。由于 $\mathcal{M}_0(\mathcal{A})$ 是单调类，我们得到
\[
\bigcup_{k=1}^\infty E_k=\bigcup_{j=1}^\infty F_j\in\mathcal{M}_0(\mathcal{A}).
\]
因此 $\mathcal{M}_0(\mathcal{A})$ 是一个 $\sigma$--代数。由于 $\mathcal{M}_0(\mathcal{A})$ 是包含 $\mathcal{A}$ 的 $\sigma$--代数，故 $\mathcal{M}_0(\mathcal{A})\supseteq\sigma_0(\mathcal{A})$。另一方面，$\sigma$--代数 $\sigma_0(\mathcal{A})$ 是包含 $\mathcal{A}$ 的单调类，故 $\sigma_0(\mathcal{A})\supseteq\mathcal{M}_0(\mathcal{A})$。于是结论成立。
\end{proof}

我们定义两个可测空间 $(X_1,\mathcal{A}_1)$ 与 $(X_2,\mathcal{A}_2)$ 的乘积。考虑可测矩形族 $\mathcal{R}\subseteq\mathcal{P}(X_1\times X_2)$，
\[
\mathcal{R}\triangleq \{E_1\times E_2\mid E_1\in\mathcal{A}_1,\ E_2\in\mathcal{A}_2\}.
\]

\begin{definition}\label{definition:ch01-乘积sigma代数}
\label{subsection:ch01-乘积}
设 $(X_1,\mathcal{A}_1)$ 与 $(X_2,\mathcal{A}_2)$ 为可测空间。最小 $\sigma$--代数 $\sigma_0(\mathcal{R})$，记为 $\mathcal{A}_1\times\mathcal{A}_2$，称为乘积 $\sigma$--代数。可测空间 $(X_1\times X_2,\mathcal{A}_1\times\mathcal{A}_2)$ 称为乘积可测空间。
\end{definition}

设 $E\subseteq X_1\times X_2$。集合
\begin{align}
E_{x_1}\triangleq \{x_2\in X_2\mid(x_1,x_2)\in E\}\qquad(x_1\in X_1) \label{align:ch01-截面}
\end{align}
称为 $E$ 的 $x_1$--截面。$x_2\in X_2$ 的 $x_2$--截面类似地定义。注意到
\[
(E^c)_{x_i}=(E_{x_i})^c,\qquad
\left(\bigcup_{n=1}^\infty E_n\right)_{x_i}=\bigcup_{n=1}^\infty(E_n)_{x_i}\qquad(i=1,2).
\]

\begin{proposition}\label{proposition:ch01-截面可测}\label{proposition:ch01-截口}\label{proposition:ch01-截面}
设 $(X_1,\mathcal{A}_1)$ 与 $(X_2,\mathcal{A}_2)$ 为可测空间，$E\in\mathcal{A}_1\times\mathcal{A}_2$。则对任意 $x_1\in X_1$ 有 $E_{x_1}\in\mathcal{A}_2$，且对任意 $x_2\in X_2$ 有 $E_{x_2}\in\mathcal{A}_1$。
\end{proposition}

\begin{proof}
令
\[
\Omega\triangleq \{E\in\mathcal{A}_1\times\mathcal{A}_2\mid E_{x_1}\in\mathcal{A}_2\ \forall x_1\in X_1\}.
\]
容易验证 $\Omega$ 是一个包含可测矩形族 $\mathcal{R}$ 的 $\sigma$--代数，故 $\Omega=\mathcal{A}_1\times\mathcal{A}_2$。类似地可证对任意 $x_2\in X_2$ 有 $E_{x_2}\in\mathcal{A}_1$，于是结论成立。
\end{proof}

\subsection{Borel $\sigma$ 代数}\label{subsection:ch01-Borel代数}

\begin{definition}\label{definition:ch01-Borel代数}
设 $(X,\mathcal{T})$ 为一拓扑空间。由拓扑 $\mathcal{T}$ 生成的 $\sigma$--代数 $\sigma_0(\mathcal{T})$（记为 $\mathcal{B}\equiv\mathcal{B}(X)\equiv\mathcal{B}(X,\mathcal{T})$）称为 Borel $\sigma$--代数。每个集合 $E\in\mathcal{B}$ 称为一个 Borel 集。
\end{definition}

\begin{example}\label{example:ch01-GdeltaFsigma}
开集的可数交集族 $\mathcal{G}_\delta$ 与闭集的可数并集族 $\mathcal{F}_\sigma$ 都包含于 $\mathcal{B}$。
\end{example}

\begin{definition}\label{definition:ch01-Baire代数}
设 $(X,\mathcal{T})$ 为一紧拓扑空间，$\mathcal{G}_\delta\supseteq G_\delta$ 为属于 $\mathcal{G}_\delta$ 的紧集族。由族 $G_\delta$ 生成的 $\sigma$--代数（记为 $\mathcal{B}_a(X)$）称为 Baire $\sigma$--代数。每个集合 $E\in\mathcal{B}_a$ 称为一个 Baire 集。
\end{definition}

\begin{remark}\label{remark:ch01-实直线Borel}
设 $X=\mathbb{R}$，$\mathcal{T}=\mathcal{T}(\mathbb{R})$ 为实拓扑。容易看出，Borel $\sigma$--代数 $\mathcal{B}(\mathbb{R})$ 与由有界开区间族生成的 $\sigma$--代数重合。它也可由其他区间族生成，例如由半直线族 $\{(a,\infty)\mid a\in\mathbb{R}\}$ 生成。类似的注记对 Borel $\sigma$--代数 $\mathcal{B}(\mathbb{R}^N)$（$N\geqslant 2$）也成立。
\end{remark}

我们提及以下关于 Borel $\sigma$--代数乘积的结果。

\begin{theorem}\label{theorem:ch01-Borel乘积}\label{theorem:ch01-Borel积}
设 $(X_1,\mathcal{T}_1)$ 与 $(X_2,\mathcal{T}_2)$ 为具有可数基的拓扑空间，$(X_1\times X_2,\mathcal{T}_1\times\mathcal{T}_2)$ 为乘积拓扑空间。则
\[
\mathcal{B}(X_1,\mathcal{T}_1)\times\mathcal{B}(X_2,\mathcal{T}_2)=\mathcal{B}(X_1\times X_2,\mathcal{T}_1\times\mathcal{T}_2).
\]
特别地，$\mathcal{B}(\mathbb{R}^M)\times\mathcal{B}(\mathbb{R}^N)=\mathcal{B}(\mathbb{R}^{M+N})$（$M,N\in\mathbb{N}$）。
\end{theorem}




\end{document}